\textbf{Phần 1. Trắc nghiệm lựa chọn}
\setcounterpageref{ex}{0}

\begin{ex} %[2D8N41]
Chia đều 8 cái kẹo cho 2 bạn. Mỗi bạn được số cái kẹo là
\choice
{$3$ cái}
{$5$ cái}
{$2$ cái}
{\True $4$ cái}
\loigiai{
$8 : 2 = 4$ cái kẹo.\\
\textbf{Đáp án của câu hỏi là D}
}
\end{ex}

\begin{ex} %[2D8N41]
Từ phép nhân $2 \times 4 = 8$, ta viết được phép chia là
\choice
{$8 : 4 = 3$}
{\True $8 : 2 = 4$ và $8 : 4 = 2$}
{$8 : 8 = 1$}
{$4 : 2 = 2$}
\loigiai{
Từ $2 \times 4 = 8$ lập được $8 : 2 = 4$ và $8 : 4 = 2$.\\
\textbf{Đáp án của câu hỏi là B}
}
\end{ex}

\begin{ex} %[2D8H41]
Từ phép nhân $5 \times 3 = 15$, phép chia tương ứng là
\choice
{\True $15 : 5 = 3$ và $15 : 3 = 5$}
{$15 : 2 = 7$}
{$15 : 4 = 3$}
{$10 : 2 = 5$}
\loigiai{
$15 : 5 = 3$ và $15 : 3 = 5$.\\
\textbf{Đáp án của câu hỏi là A}
}
\end{ex}

\begin{ex} %[2D8H41]
Kết quả của phép chia $10 : 2$ là
\choice
{$4$}
{$6$}
{$8$}
{\True $5$}
\loigiai{
$10 : 2 = 5$.\\
\textbf{Đáp án của câu hỏi là D}
}
\end{ex}

\begin{ex} %[2D8H41]
Kết quả của phép chia $20 : 5$ là
\choice
{$3$}
{$5$}
{\True $4$}
{$6$}
\loigiai{
$20 : 5 = 4$.\\
\textbf{Đáp án của câu hỏi là C}
}
\end{ex}

\begin{ex} %[2D8H41]
Có 12 quả cam chia đều vào 2 đĩa. Mỗi đĩa có số quả cam là
\choice
{\True $6$ quả}
{$5$ quả}
{$7$ quả}
{$8$ quả}
\loigiai{
$12 : 2 = 6$ quả cam.\\
\textbf{Đáp án của câu hỏi là A}
}
\end{ex}

\begin{ex} %[2D8V41]
Có 15 bông hoa cắm đều vào 5 lọ. Mỗi lọ có số bông hoa là
\choice
{$2$ bông}
{\True $3$ bông}
{$4$ bông}
{$5$ bông}
\loigiai{
$15 : 5 = 3$ bông hoa.\\
\textbf{Đáp án của câu hỏi là B}
}
\end{ex}

\begin{ex} %[2D8V41]
Có 18 bạn xếp thành 2 hàng đều nhau. Mỗi hàng có số bạn là
\choice
{$8$ bạn}
{$10$ bạn}
{\True $9$ bạn}
{$7$ bạn}
\loigiai{
$18 : 2 = 9$ bạn.\\
\textbf{Đáp án của câu hỏi là C}
}
\end{ex}

\begin{ex} %[2D8V41]
Có 25 quyển vở chia đều cho 5 bạn. Mỗi bạn nhận được số quyển vở là
\choice
{\True $5$ quyển}
{$4$ quyển}
{$6$ quyển}
{$7$ quyển}
\loigiai{
$25 : 5 = 5$ quyển vở.\\
\textbf{Đáp án của câu hỏi là A}
}
\end{ex}

\begin{ex} %[2D8V41]
Điền dấu thích hợp: $14 : 2 \dots 30 : 5$.
\choice
{$<$}
{$=$}
{\True $>$}
{Không so sánh được}
\loigiai{
$14 : 2 = 7$ và $30 : 5 = 6$. Vì $7 > 6$ nên điền dấu $>$.\\
\textbf{Đáp án của câu hỏi là C}
}
\end{ex}

\begin{ex} %[2D8V41]
Tính giá trị biểu thức: $16 : 2 + 25$.
\choice
{$31$}
{$35$}
{$37$}
{\True $33$}
\loigiai{
$8 + 25 = 33$.\\
\textbf{Đáp án của câu hỏi là D}
}
\end{ex}

\begin{ex} %[2D8C41]
Có 35 cái kẹo chia đều vào 5 túi, lấy ra 2 túi kẹo. Số kẹo lấy ra là
\choice
{$12$ cái}
{\True $14$ cái}
{$15$ cái}
{$16$ cái}
\loigiai{
Mỗi túi có: $35 : 5 = 7$ cái kẹo.\\
2 túi có: $7 \times 2 = 14$ cái kẹo.\\
\textbf{Đáp án của câu hỏi là B}
}
\end{ex}

\textbf{Phần 2. Trắc nghiệm đúng sai}
\setcounterpageref{ex}{0}

\begin{ex} %[2D8H41]
Xét mối quan hệ giữa phép nhân và phép chia:
\choiceTF
{\True Từ phép nhân $2 \times 6 = 12$ ta có $12 : 2 = 6$.}
{\True Từ phép nhân $2 \times 6 = 12$ ta có $12 : 6 = 2$.}
{$12 : 2 = 5$.}
{\True $12 : 2 = 6$.}
\loigiai{
\begin{itemize}
    \item $12 : 2 = 6$: Đúng.
    \item $12 : 6 = 2$: Đúng.
    \item $12 : 2 = 6 \neq 5$. Ý 3 Sai, Ý 4 Đúng.
\end{itemize}
\begin{itemize}
    \item a) \textbf{Đúng}.
    \item b) \textbf{Đúng}.
    \item c) \textbf{Sai}.
    \item d) \textbf{Đúng}.
\end{itemize}
}
\end{ex}

\begin{ex} %[2D8V41]
Có 20 cái bánh chia đều vào các hộp:
\choiceTF
{\True Nếu chia vào 2 hộp thì mỗi hộp có 10 cái bánh ($20 : 2 = 10$).}
{\True Nếu chia vào 5 hộp thì mỗi hộp có 4 cái bánh ($20 : 5 = 4$).}
{Nếu chia vào 4 hộp thì mỗi hộp có 6 cái bánh.}
{\True Nếu chia vào 4 hộp thì mỗi hộp có $20 : 4 = 5$ cái bánh.}
\loigiai{
\begin{itemize}
    \item 2 hộp: $20 : 2 = 10$: Đúng.
    \item 5 hộp: $20 : 5 = 4$: Đúng.
    \item 4 hộp: $20 : 4 = 5 \neq 6$. Ý 3 Sai, Ý 4 Đúng.
\end{itemize}
\begin{itemize}
    \item a) \textbf{Đúng}.
    \item b) \textbf{Đúng}.
    \item c) \textbf{Sai}.
    \item d) \textbf{Đúng}.
\end{itemize}
}
\end{ex}

\textbf{Phần 3. Trả lời ngắn}
\setcounterpageref{ex}{0}

\begin{ex} %[2D8H41]
Tính: $14 : 2$.
\shortans{7}
\loigiai{
\textbf{Đáp số:} $7$\\
$14 : 2 = 7$.
}
\end{ex}

\begin{ex} %[2D8H41]
Tính: $35 : 5$.
\shortans{7}
\loigiai{
\textbf{Đáp số:} $7$\\
$35 : 5 = 7$.
}
\end{ex}

\begin{ex} %[2D8V41]
Có 16 lít dầu rót đều vào 2 can. Hỏi mỗi can có bao nhiêu lít dầu?
\shortans{8}
\loigiai{
\textbf{Đáp số:} $8$\\
$16 : 2 = 8\text{ l}$.
}
\end{ex}

\begin{ex} %[2D8V41]
Có 40 cái kẹo chia đều cho 5 bạn. Hỏi mỗi bạn được bao nhiêu cái kẹo?
\shortans{8}
\loigiai{
\textbf{Đáp số:} $8$\\
$40 : 5 = 8$ cái.
}
\end{ex}

\begin{ex} %[2D8V41]
Tính giá trị biểu thức: $45 : 5 + 16$.
\shortans{25}
\loigiai{
\textbf{Đáp số:} $25$\\
$9 + 16 = 25$.
}
\end{ex}

\begin{ex} %[2D8C41]
Một khúc gỗ dài $18\text{ dm}$ được cưa thành các đoạn bằng nhau dài $2\text{ dm}$. Hỏi cần bao nhiêu lần cưa?
\shortans{8}
\loigiai{
\textbf{Đáp số:} $8$\\
Số đoạn gỗ: $18 : 2 = 9$ đoạn.\\
Số lần cưa: $9 - 1 = 8$ lần.
}
\end{ex}
